\section{Algorithm and numerical implementation}\label{sec:algorithm}
The implementation exposes two cut policies (ESH and ECP), two master formulations (hull and big-$M$), and two execution modes (single-tree and repeated-master, termed multi-tree). These choices form the matched experimental comparisons. The following description concerns the frozen implementation used in the study. Its numerical safeguards and finite stopping rules are distinguished from the sufficient exact-arithmetic contracts of \cref{sec:theory}.

\subsection{Initialization and the common separation routine}
The extractor retains the original GDP rows and their term membership, converts the objective to minimization, and introduces an epigraph for a nonlinear objective. Affine bound propagation uses global rows and unions of branchwise bound deductions. LP optimization over global affine rows attempts to recover missing bounds required by a formulation. These are floating-point preprocessing operations. The master refuses a required unbounded hull copy or a big-$M$ coefficient that cannot be calculated finitely. In particular, a nominal default big-$M$ option is not used as a substitute for a missing valid row bound in the reviewed construction.

Both policies run the same interior initialization. For each term's nonlinear row collection, and separately for global nonepigraph nonlinear rows, an auxiliary NLP minimizes a common slack $s$ subject to $g_j(x)\le s$, relevant affine rows, and variable bounds. Original integer coordinates are continuous in this auxiliary problem. Otherwise unbounded coordinates receive an artificial finite search box around their starting values (scale $10^3$), and the common slack has lower bound $-10^6$. These restrictions aid anchor search only; they do not become original-model or master bounds and do not establish the compactness premise of the theory. The implementation tries up to four starts, with deterministic restarts, and retains a candidate anchor if the loaded slack is below $-10^{-6}$. This is only a candidate for radial use: the separator checks the actual value of the particular row at that anchor, requiring it to be below $-10^{-9}$. Failure to obtain a usable anchor invokes ECP and never proves term infeasibility. The epigraph row uses ECP because it is excluded from the anchor NLP. Keeping interior initialization in ECP isolates the subsequent cut-policy change but imposes work an independently designed ECP solver need not perform.

Initial point tangents are added for every nonlinear row, using anchor coordinates when available and otherwise bounded starting values. Invalid or nonfinite evaluations trigger a bounded deterministic search for a finite starting tangent (at most 31 attempts per row), followed by an error if none is found. All later tangents retain the full affine constant in \eqref{eq:tangent}.

\begin{quote}
\textbf{Common separation routine at $(x,y,\nu)$ and cutoff $\tau$.}
\begin{enumerate}
 \item Test global nonlinear rows at $(x,y)$. For each term with $y_{ik}>\tau$, test its rows at $p_{ik}=\nu_{ik}/y_{ik}$ in the hull master, or at $x$ in the big-$M$ master.
 \item For every tested row with value exceeding $\min\{\texttt{cut\_tol},\texttt{feas\_tol}\}$, use its point tangent for ECP. For ESH, if an anchor contains all row variables and its row value is below $-10^{-9}$, bisect that row along the anchor--candidate segment. Stop after at most 60 bisections or once the parameter bracket has width below $10^{-10}$. Linearize at the exterior bracket endpoint.
 \item Use the radial tangent if its coefficients are finite and its untransformed violation at the test point exceeds $0.1\,\texttt{cut\_tol}$; otherwise use a point tangent. Nonfinite candidate row values or fallback coefficients raise an evaluation error.
 \item Transform each term tangent using \eqref{eq:perspective-cut} or \eqref{eq:bigm-cut}, and submit it as an ordinary master constraint, an integer-candidate lazy constraint, or a fractional user cut according to the caller.
\end{enumerate}
\end{quote}
The line search is rowwise, rather than a single root search on the maximum of all rows. The exact boundary target belongs to the boundary of that row's sublevel set and need not belong to the boundary, or even the feasible set, of the full multirow disjunct. The implemented bisection instead linearizes at an approximate exterior bracket endpoint. Term structure determines the row tags, anchors and cut transformation; it does not turn the line search into projection onto the whole term. No optimization problem is solved within this cut routine. The default values of \texttt{cut\_tol} and \texttt{feas\_tol} are both $10^{-6}$. Checking the untransformed radial violation is weaker than the certified transformed-margin requirement in \cref{sec:arithmetic}, particularly at small weights. The fallback checks coefficient finiteness but supplies no certified rounding-error bound. Anchor values, intermediate bisection values, and the scalar radial tangent violation have no separate finiteness guard; returned nonfinite values therefore pass through ordinary comparisons. Thrown evaluation exceptions propagate. Thus this routine implements a practical numerical policy, not the certified scheme in \cref{prop:arithmetic}.

Hull normalization also has a domain qualification. In exact arithmetic the scaled bounds imply $\nu_{ik}/y_{ik}\in C$. Numerically, a scaled-bound error can be amplified by division by a small positive weight. The frozen routine neither projects the normalized point onto $C$ nor separately verifies its domain membership. Finite row values and derivatives at a point outside the assumed convexity domain do not establish a valid supporting inequality. The exact validity results require in-domain evaluation points; this numerical premise is not certified by the routine. This is a potential numerical risk, not an observed failure of the reported experiments.

\subsection{LP initialization and the two tree policies}
\begin{quote}
\textbf{Common LP phase.}
\begin{enumerate}
 \item Relax all master integer variables, including original integer coordinates and GDP indicators. Solve the current LP with the remaining time budget.
 \item When the LP is optimal, retain the best attained objective as the initial LP bound. At any usable LP solution, record the maximum numerical weighted perspective residual over \emph{all} positive-weight hull blocks, including those below the separation cutoff.
 \item Call the separation routine with its default $\tau=10^{-6}$. Repeat after adding cuts unless the routine adds none, a resource/status limit occurs, or the objective stagnates.
 \item Restore integrality and proceed to the selected mixed-integer execution mode.
\end{enumerate}
\end{quote}
The default LP cap is 200 iterations. The recorded \texttt{lp\_iters} counts iterations of this initial LP separation phase, not internal simplex or barrier iterations or all node LP solves. One phase iteration can include a second solve with dual reductions disabled after an ambiguous or unbounded status. Stagnation is detected after at least six cut-producing solves when the newest objective differs from the one five iterations earlier by at most $10^{-5}(1+|v_{\rm LP}|)$. The recorded exit distinguishes no separating cuts, stagnation, iteration limit, time limit and master status. A no-cut exit is only a completed test of the rows visited under the numerical cutoff; a stagnation or limit exit does not provide even that assertion. The residual diagnostic is absent for big-$M$ and is reported as infinite if a required evaluation is undefined. The retained LP bound is still the best optimal LP value reached even if cuts were added after its solve. It is a numerical solver bound, not an interval certificate.

\begin{quote}
\textbf{Multi-tree mode.}
\begin{enumerate}
 \item Solve the current MILP with the remaining time budget. Read its objective bound as the minimization lower bound $B$; do not use the master incumbent objective for that purpose after an incomplete solve.
 \item Separate its integer candidate using the common routine. If integer-point NLP recovery is enabled, fix its integer assignment and solve the selected continuous subproblem. Validate any recovered point, update the incumbent $U$ when it improves, and add tangents at the recovered point for global and selected-term rows.
 \item If separation added no cuts, independently validate the master point as a possible original-model incumbent. A relaxed epigraph objective alone cannot establish that incumbent.
 \item Stop when the numerical gap criterion holds or a resource/status condition occurs. If no cut was added and no incumbent exists, report a stall. Otherwise repeat with the accumulated master.
\end{enumerate}
\end{quote}
The default cap is 500 MILP iterations. Previously added ordinary cuts remain in the master. A merely tolerance-feasible incumbent is not used to add an objective cutoff in this mode. Failure to recover a point from an NLP does not create an infeasibility cut or a no-good cut for that assignment.

\begin{quote}
\textbf{Single-tree mode.}
\begin{enumerate}
 \item Start one MILP branch-and-cut solve with lazy constraints enabled and the remaining time budget.
 \item At every integer-solution callback, run separation and submit its cuts as lazy constraints. If no cut is generated, independently validate the original point before recording an incumbent.
 \item When enabled, perform integer-point NLP recovery at that callback. Validate an improving result and submit it as a heuristic solution, setting hull copies to the recovered $x$ in selected terms and zero in inactive terms.
 \item If optional fractional user cuts are enabled, separate optimal node relaxations at the root and at node counts divisible by 50, using $\tau=0.05$. Submit those cuts as user cuts.
 \item On solver return, validate any final master incumbent, obtain the master bound, and apply the numerical incumbent/gap/status rules below.
\end{enumerate}
\end{quote}
The main study enables integer-point NLP recovery at every integer callback and disables fractional user cuts. The source also offers an adaptive recovery option, but that option is not the main experimental policy. NLP-solution tangents are explicitly added in multi-tree mode; the single-tree NLP path instead submits the validated heuristic point and relies on callback separation. Optional fractional cuts are an ablation, use a larger fixed cutoff than the LP phase, and have no explicit finite total cut cap. The implementation does not add an external mechanism that verifies old-lazy-cut satisfaction before generating fresh cuts. The corresponding solver properties in \cref{thm:single-tree} are therefore assumptions for the theoretical algorithm, not proven properties of this callback implementation.

\subsection{NLP recovery, validation, and failure states}
The recovery NLP fixes \emph{all} original integer variables and indicators to the current rounded assignment, deactivates inactive disjuncts, and solves the continuous selected model on a clone. Fixed original variables are preserved when preparing its start. Validated results are cached by integer assignment; repeated failures can also be cached to avoid indefinite retries. An NLP return code is never sufficient for incumbent acceptance. A failed or locally infeasible NLP return is not a proof that the assignment has no feasible point.

Before acceptance, the internal validator requires finite original variable values, normalizes integer values only when they are within \texttt{feas\_tol} of an integer, and restores fixed values only within that same tolerance. It then checks variable bounds, global affine and nonlinear rows, exclusive selection and selected-term rows on this normalized point. Indicator logic is included among the global affine rows. It reevaluates the original objective on that point and resets any artificial epigraph to the resulting minimization value. Inactive-term constraints are not required. The stored incumbent and final model writeback use the same normalized values. A feasible point need not be an optimal NLP solution to improve the incumbent.

For the main settings, internal acceptance allows maximum absolute original-row and bound violation $10^{-6}$. With internal minimization incumbent $U$ and bound $B$, the implementation reports a closed numerical gap only when both are finite and
\begin{equation}\label{eq:internal-gap}
 |U-B|\le 10^{-6}+10^{-4}|U|.
\end{equation}
The symmetric test permits small numerical bound inconsistency and rejects large inconsistency; it is not the exact upper/lower-bound inequality of \cref{thm:single-tree}. The benchmark independently rechecks the original witness with a scale-dependent feasibility tolerance and a separate gap formula, specified in the experimental design below. Consequently raw solver status and an accepted experimental solve are separate quantities.

An internal \texttt{optimal} status means a numerically accepted incumbent and a closed numerical gap under \eqref{eq:internal-gap}. Results explicitly record \texttt{rigorous\_certificate=False}, the feasibility tolerance, and the maximum incumbent residual. Time limits, iteration limits, stalls, cutoffs, and other master statuses remain distinguishable. Callback exceptions terminate the solve with \texttt{numerical\_error} and preserve the exception detail. Evaluation errors outside callbacks can propagate as failed runs. A solver infeasibility report after an accepted incumbent is treated as a numerical inconsistency, not as a successful infeasibility proof. Ambiguous master termination is not silently relabeled infeasible.

\subsection{Supported scope and accounting}
The reviewed input scope consists of flat exclusive disjunctions with convex differentiable inequality rows and adequate finite bounds for the selected formulation. Unsupported structures are refused: nonexclusive disjunctions, nested or multiply assigned disjuncts, active orphan disjuncts, objectives inside disjuncts, SOS constraints, nonlinear equalities, two-sided nonlinear rows, and nonfixed native GDP indicators inside disjunct bodies. Constant contradictory rows are retained so that an infeasible term can be excluded without excluding its alternatives. Declared variables remain subject to their bounds even if absent from all expressions. Successful extraction does not certify convexity, smoothness or safe evaluation throughout the box; these remain modeling obligations. The external-model analysis therefore distinguishes smooth supported expressions from norm and domain-sensitive cases.

The time budget includes construction and preprocessing. Remaining time is passed to MILP/LP solves and, for the default Ipopt interface, as an NLP CPU-time limit. These limits do not guarantee hard wall-clock preemption; the experimental process wrapper supplies a separate wall limit. Interior initialization, reduced NLP work, row separation and total wall time are recorded separately. The \emph{recorded row-separation time per run}, stored as \texttt{time\_cuts}, accumulates elapsed time in completed calls to \texttt{\_esh\_cuts}: row checks, radial searches when used, and linearization and coefficient work. It excludes initial tangents, NLP-solution tangents, normalization and transformation work outside that routine, and installation of cuts in the master or callbacks. The cut count has a broader scope: it also includes initial and NLP-solution tangents. In single-tree mode the MILP solve timer includes callbacks and therefore overlaps NLP and separation timers. Such times cannot be summed as disjoint costs. NLP recovery disabled means no integer-point recovery NLPs; the common interior NLP initialization remains. This distinction is essential for interpreting that ablation.
